\section{No Longer Servants}

John 15 places commandment, obedience, love, and friendship in one sequence. Christ says, ``If you keep my commandments, you will abide in my love''; then, ``This is my commandment, that you love one another as I have loved you''; and finally, ``I have called you friends.''\endnote{John 15:9--17. Abiding, commandment, Christ's self-gift, friendship, and mission interpret one another.} Friendship does not abolish obedience. It changes the relation in which obedience is rendered.

Irenaeus reads the same conjunction without relaxing either term. Christ does not abolish the natural precepts but extends them from external restraint to the desires of the heart. The liberated person remains subject to God; indeed, Irenaeus says the work of liberty is greater and more glorious than servile obedience. Servants and children owe piety to the same Lord, but children follow with confidence, and friends know what the Word has received from the Father.\endnote{Irenaeus, \emph{Against Heresies} IV.13.1--4 and IV.14.1: friendship perfects the creature, not God.}

The transition benefits the creature, not God. Irenaeus immediately insists that God did not accept Abraham's friendship out of need and does not profit from human service. The perfect God grants life to those who follow him; obedience perfects the servant rather than completing the Lord (IV.13.4; IV.14.1). Creaturely response remains asymmetric even at its most intimate.

Aquinas supplies friendship with a precise grammar. Not every love is friendship. Friendship requires willing the good of another, reciprocal love, and some communication or shared life. God communicates to creatures a participation in eternal beatitude; upon that gift a real friendship with God can be founded. This friendship is charity.\endnote{ST II--II, q.~23, a.~1. The shared good is God's prior gift; friendship remains real and unequal.}

No equality of nature follows. The communication is first given by God, charity in the soul is created, and the final good remains possessed by God as source and by the creature through participation (q.~23, a.~2; q.~24, a.~2). Christ remains Lord when he calls the disciples friends. The Creator--creature distinction is not relaxed so that communion can begin. Communion is possible because divine transcendence is not rivalry and grace can perfect a creature without making it uncreated.

What changes is the creature's manner of willing. Fear becomes servile when punishment is treated as the chief evil; mercenary love calculates the wage. Filial fear dreads separation from the beloved; charity cleaves to divine goodness itself. Augustine says punitive fear recedes as charity grows, while chaste fear remains watchful not to lose communion with the beloved. Aquinas gives the distinction formal precision: charity casts out servility; fear of punishment may remain but diminishes as charity grows; filial fear increases.\endnote{Augustine, \emph{Homilies on First John} IX.4--6; ST II--II, q.~19, aa.~10--11; CCC 1828, 1972.}

Fear of punishment can belong to the beginning of conversion and may remain as charity grows, though servility recedes. The theological virtue itself is not measured by sensible closeness to God. Friendship names an objective participation and a stable orientation of will, lived through acts that may be costly, obscure, and dry.

Charity does not remove the command. It begins to remove the distance between the will and the good that the command protects. The friend obeys because, by grace, that good is already beginning to be loved.
